\content{
\begin{example} Find the equation of the plane which passes through the point
$(2,4,-1)$ with normal vector $\vec{n}=\la 2,3,4\ra $.
\end{example}
}{% presentation
\vspace{2in}
}{% nopresentation
\vspace{2in}
}{% comments
To sketch this plane, find the intercepts at the axes and draw a triangle.
}

\content{
\begin{example} Find the equation of the plane which passes through the points 
$P(1,3,2)$, $Q(3,-1,6)$ and $R(5,2,0)$.
\end{example}
}{% presentation
\vspace{3in}
}{% nopresentation
\vspace{2in}
}{% comments
}


\content{
\begin{example}
Find the point at which the line $x = 2+3t$, $y=-4t$, $z=5+t$ intercepts the
plane $4x+5y-2z=18$.
\end{example} 
}{% presentation
\vspace{2in}
}{% nopresentation
\vspace{1in}
}{% comments
Simply insert the line equations into plane equation and solve for $t$.
}

\content{
\begin{example} Find the angle between the two planes $x+y+z=1$ and $x-2y+3z=1$. 
\end{example}
}{% presentation
\vspace{2in}
}{% nopresentation
\vspace{2in}
}{% comments
Note that the angle between two planes is defined to be the angle between their
normal vectors. 

Answer is $\cos^{-1}\left(\frac{2}{\sqrt{42}}\right) \approx 72^\circ$.
}

\subsubsection{Distances}

\content{
\vspace{3in}
}{% presentation
}{% nopresentation
}{% comments
Show that the distance from a point $P$ to a plane $D$ is given by taking any
point $P_0$ in the plane, and computing the length of the projection of the
vector $\vec{P_0P}$ onto the normal vector of $D$.
}

\content{
\begin{example} Find the distance from the point $(4,0,1)$ to the plane
$x-y+z=1$.
\end{example}
}{% presentation
\vspace{3in}
}{% nopresentation
\vspace{2in}
}{% comments
}

\content{
\begin{example} Find the distance between the two planes $10x+2y-2z=5$ and 
$5x+y-z=1$. 
\end{example}
}{% presentation
\vspace{3in}
}{% nopresentation
\vspace{2in}
\newpage
}{% comments
Notice that these two planes are parallel. So we can take any point on one plane
and compute it's distance to the other plane. 
}

\subsection{Cylinders and Quadric Surfaces}

\content{
\begin{definition} A cylinder is a surface that consists of all lines that are
parllel to a given line and pass through a given plane curve. 
\end{definition}
}{% presentation
\vspace{4in}
}{% nopresentation
\vspace{4in}
}{% comments
Give several examples, including a regular cylinder, and also some other
cylinders. 

Note that these are infinitly tall, or long. 
}

\content{
\begin{definition} A quadric surface is the graph of a second-degree equation in three
variables. 
\end{definition}

\begin{example} Use traces to sketch the quadric surface with equation 
$$ x^2 + \frac{y^2}{9}+\frac{z^2}{4}=1. $$
\end{example}
}{% presentation
\vspace{4in}
}{% nopresentation
\vspace{4in}
}{% comments
Level curves are elipses, as long as $|z|\leq 4$.

This surface is known as an ellipsoid. 
}

\content{
\begin{example} Sketch the surface 
$$ z=4x^2+y^2.$$
\end{example}
}{% presentation
\vspace{4in}
}{% nopresentation
\vspace{4in}
}{% comments
This surface is known as an elliptic paraboloid. 
}

\content{
\begin{example} Sketch the surface
$$ \frac{x^2}{4}+y^2-\frac{z^2}{4} = 1. $$
\end{example}
}{% presentation
\vspace{4in}
}{% nopresentation
\vspace{4in}
}{% comments
This surface is a hyperboloid of one sheet. 

\textbf{There is a table of quadric surfaces to see in the book.}
}

